\documentclass[11pt,a4paper]{article}

\usepackage{iftex}
\ifPDFTeX
  \usepackage[utf8]{inputenc}
  \usepackage[T1]{fontenc}
  \usepackage{lmodern}
\else
  \usepackage{fontspec}
\fi
\usepackage[spanish]{babel}
\usepackage[margin=2.5cm]{geometry}
\usepackage{booktabs}
\usepackage{array}
\usepackage{graphicx}
\usepackage{enumitem}
\usepackage{hyperref}

\hypersetup{
  colorlinks=true,
  linkcolor=black,
  urlcolor=black,
  pageanchor=false,
  pdftitle={Laboratorio 6 - Visualización de Datos con Power BI},
  pdfauthor={Torres Lezama, Mathias James}
}
\setlength{\parindent}{0pt}
\setlength{\parskip}{0.6em}

\newcommand{\pregunta}[2]{\subsection*{Pregunta #1}\textit{#2}}

\begin{document}

% ===================================================================
% CARATULA
% ===================================================================
\begin{titlepage}
\thispagestyle{empty}
\begin{center}

{\LARGE \textbf{UNIVERSIDAD NACIONAL MAYOR DE SAN MARCOS}}\\[3pt]
{\itshape\small Decana de América, fundada en 1551}\\[40pt]

% El logo es opcional: si el archivo no existe en img/, el informe compila igual.
\IfFileExists{img/logo_unmsm_0.png}{\includegraphics[width=3.2cm]{img/logo_unmsm_0.png}\\[40pt]}{\vspace{60pt}}

{\large \textbf{FACULTAD DE INGENIERÍA DE SISTEMA E INFORMÁTICA}}\\[6pt]
{\large \textbf{ESCUELA PROFESIONAL DE INGENIERÍA DE SOFTWARE}}\\[80pt]

{\Large \textbf{Laboratorio 6 - Visualización de Datos con Power BI}}\\[90pt]

\end{center}

\begin{center}
    \large
    \textbf{Estudiante:}\\
    Torres Lezama, Mathias James\\[20pt]

    \textbf{Docente:}\\
    Gamarra Moreno, Juan\\[20pt]

    \textbf{Asignatura:}\\
    Inteligencia de Negocios\\
\end{center}

\vfill

\begin{center}
\large Lima, Perú\\
2026
\end{center}

\end{titlepage}

\section*{Actividad propuesta: Dashboard de Clientes Bancarios}

Este documento recoge las respuestas de las preguntas de análisis, elaboradas a partir del archivo \texttt{clientes\_banco\_retail.csv} (400 clientes) y del informe \texttt{Lab06\_TorresLezamaMathias.pbix}.

\section*{Resumen general del portafolio}

Los indicadores globales del dashboard son los siguientes:

\begin{center}
\begin{tabular}{@{}lr@{}}
\toprule
Indicador & Valor \\
\midrule
Total de clientes & 400 \\
Clientes en mora & 76 \\
Tasa de mora (\%) & 19.00 \\
Saldo promedio total (S/) & 40,617.10 \\
Satisfacción promedio (escala 1 a 5) & 3.01 \\
\bottomrule
\end{tabular}
\end{center}

\section*{Preguntas de análisis}

% ---------------------------------------------------------------
\pregunta{1}{¿Cuál es el producto financiero con la mayor cantidad de clientes en mora? ¿A qué segmento etario pertenecen principalmente estos clientes?}

El producto con más clientes en mora es \textbf{Depósito Plazo Fijo}, con 23 clientes, seguido de Cuenta Ahorros (17) y Tarjeta Crédito (13).

\begin{center}
\begin{tabular}{@{}lr@{}}
\toprule
Producto & Clientes en mora \\
\midrule
Depósito Plazo Fijo & 23 \\
Cuenta Ahorros & 17 \\
Tarjeta Crédito & 13 \\
Crédito Personal & 10 \\
Cuenta Corriente & 8 \\
Hipotecario & 5 \\
\bottomrule
\end{tabular}
\end{center}

Dentro de Depósito Plazo Fijo, el segmento con más clientes en mora es \textbf{Adulto (31-50)}, con 9 clientes. Sin embargo, la diferencia es pequeña: Senior (51+) tiene 8 y Joven (18-30) tiene 6. Por ello, la mora en este producto no se concentra de forma marcada en un solo segmento.

% ---------------------------------------------------------------
\pregunta{2}{¿En qué departamento del Perú el banco tiene el mayor saldo promedio? ¿Coincide con el departamento de mayor número de clientes?}

El departamento con mayor saldo promedio es \textbf{Lambayeque}, con S/ 46,863.64. \textbf{No coincide} con el departamento de mayor número de clientes, que es \textbf{Lima}, con 158 clientes y un saldo promedio de S/ 39,661.57.

\begin{center}
\begin{tabular}{@{}lrr@{}}
\toprule
Departamento & Clientes & Saldo promedio (S/) \\
\midrule
Lambayeque & 14 & 46,863.64 \\
Junín & 27 & 45,860.25 \\
Ica & 17 & 45,276.27 \\
Piura & 28 & 43,000.09 \\
Cajamarca & 21 & 40,381.93 \\
La Libertad & 31 & 40,143.33 \\
Lima & 158 & 39,661.57 \\
Cusco & 33 & 39,240.02 \\
Áncash & 13 & 38,724.87 \\
Arequipa & 58 & 38,301.48 \\
\bottomrule
\end{tabular}
\end{center}

Lima concentra el volumen de clientes, pero su saldo promedio está por debajo del de Lambayeque, Junín, Ica y Piura. El mayor saldo promedio de Lambayeque se apoya en una base pequeña (14 clientes), por lo que el promedio es más sensible a valores extremos.

% ---------------------------------------------------------------
\pregunta{3}{¿Cuál es el canal de preferencia más utilizado por los clientes del segmento <<Senior (51+)>>? ¿Qué implicaciones tendría este dato para la estrategia digital del banco?}

El canal más utilizado por los clientes Senior es \textbf{Call Center}, con 28 de 96 clientes (29.2\%), seguido de Cajero Automático (21). El canal menos utilizado es App Móvil (14).

\begin{center}
\begin{tabular}{@{}lrr@{}}
\toprule
Canal preferido & Clientes Senior & Porcentaje \\
\midrule
Call Center & 28 & 29.2\% \\
Cajero Automático & 21 & 21.9\% \\
Web & 18 & 18.8\% \\
Agencia & 15 & 15.6\% \\
App Móvil & 14 & 14.6\% \\
\bottomrule
\end{tabular}
\end{center}

\textbf{Implicaciones para la estrategia digital:}
\begin{itemize}[nosep]
  \item Este segmento prefiere canales asistidos (telefónico y cajero) y utiliza poco la aplicación móvil, por lo que la adopción digital es baja.
  \item El banco debería diseñar programas de educación financiera digital y acompañamiento para migrar gradualmente a estos clientes a la App Móvil y la Web.
  \item Conviene mantener el Call Center como canal de soporte, con asistencia guiada para operaciones digitales, en lugar de forzar un cambio abrupto.
  \item Las interfaces para este segmento deben priorizar la simplicidad y la accesibilidad.
\end{itemize}

% ---------------------------------------------------------------
\pregunta{4}{Observe el gráfico de columnas apiladas al 100\% de mora por segmento etario. ¿Qué segmento presenta la mayor tasa de mora? ¿Qué acciones recomendaría al área de riesgos?}

El segmento con la mayor tasa de mora es \textbf{Senior (51+)}, con 22.92\%.

\begin{center}
\begin{tabular}{@{}lr@{}}
\toprule
Segmento etario & Tasa de mora (\%) \\
\midrule
Senior (51+) & 22.92 \\
Adulto (31-50) & 18.50 \\
Joven (18-30) & 16.79 \\
\bottomrule
\end{tabular}
\end{center}

\textbf{Acciones recomendadas para el área de riesgos:}
\begin{itemize}[nosep]
  \item Implementar cobranza preventiva y alertas tempranas para clientes Senior antes del vencimiento.
  \item Ofrecer opciones de renegociación o reprogramación de cuotas adaptadas a este segmento.
  \item Revisar la política de otorgamiento de crédito para Senior (ingresos, capacidad de pago y plazos).
  \item Usar canales de contacto que este segmento prefiere (Call Center) para el seguimiento de cobranza.
  \item Monitorear la evolución de la tasa de mora por segmento de forma periódica.
\end{itemize}

% ---------------------------------------------------------------
\pregunta{5}{¿Cuál es la satisfacción promedio de los clientes que están en mora versus los que están al día?}

Al filtrar con el segmento \texttt{en\_mora}, se obtiene:

\begin{center}
\begin{tabular}{@{}lr@{}}
\toprule
Estado & Satisfacción promedio (1 a 5) \\
\midrule
Al día (\texttt{en\_mora} = 0) & 2.99 \\
En mora (\texttt{en\_mora} = 1) & 3.08 \\
\bottomrule
\end{tabular}
\end{center}

Los clientes en mora presentan una satisfacción ligeramente \textbf{mayor} (3.08) que los clientes al día (2.99). El resultado es contraintuitivo, pero la diferencia es de solo 0.09 puntos en una escala de 1 a 5, por lo que no es concluyente: con esta muestra no se puede afirmar que la mora esté asociada a una mayor satisfacción.

\end{document}
